\section*{الفصل \ref{ch:g11:titration} --- المعايرة}

\begin{solution}{exo:g11:titration:1}
\omterm{def:g11:titration:titration}{المحلول المعايِر} هو \omterm{def:g2:mixing-and-dissolving:dissolve}{المحلول} المعروف التركيز بدقة، في
السحاحة؛ \omterm{def:g11:titration:titration}{والمحلول المراد معايرته} يحوي النوع المراد قياسه، في
الدورق. وعند \omterm{def:g11:titration:equivalence}{التكافؤ} يكون \omterm{def:g11:titration:titration}{المحلول المعايِر} المضاف والنوع المراد معايرته
بنسب المعادلة: فيُستهلكان كلاهما.
\end{solution}

\begin{solution}{exo:g11:titration:2}
$n = 0.0200 \times 0.0125 = \qty{2.50e-4}{mol}$.
\end{solution}

\begin{solution}{exo:g11:titration:3}
$V_E = \qty{14.3}{mL}$؛ ويبدأ خط الحديد (II) عند
\qty{14.3e-4}{mol}، أي \qty{1.43e-3}{mol}.
\end{solution}

\begin{solution}{exo:g11:titration:4}
\omterm{def:g11:titration:titration}{المحلول المعايِر} نفسه ملوّن (بنفسجي) وناتجه يكاد يكون عديم
اللون: فيظهر اللون في الدورق ما إن تصير البرمنغنات
بفائض.
\end{solution}

\begin{solution}{exo:g11:titration:5}
5: فالمعادلة تستهلك 5 \omterm{def:g9:ions:ion}{أيونات} حديد (II) لكل \omterm{def:g9:ions:ion}{أيون} برمنغنات، وعند
\omterm{def:g11:titration:equivalence}{التكافؤ} تكون \omterm{def:g7:chemical-reactions:reactant}{المتفاعلات} قد جُمعت بهذه
النسب.
\end{solution}

\begin{solution}{exo:g11:titration:6}
$n(\ce{MnO4^-}) = 0.0200 \times 0.00840 = \qty{1.68e-4}{mol}$؛
$n(\ce{Fe^{2+}}) = 5 \times \num{1.68e-4} = \qty{8.40e-4}{mol}$؛
$c = \num{8.40e-4} / 0.0100 = \qty{0.0840}{mol/L}$؛
$C_m = 0.0840 \times 55.8 = \qty{4.69}{g/L}$.
\end{solution}

\begin{solution}{exo:g11:titration:7}
النسبة 1:1: $n = 0.0100 \times 0.0114 = \qty{1.14e-4}{mol}$ في
\qty{10.0}{mL}، إذن \qty{1.14e-3}{mol} في \qty{100.0}{mL}، أي
$\num{1.14e-3} \times 176.0 = \qty{0.201}{g}$: نحو \qty{200}{mg}.
\end{solution}

\begin{solution}{exo:g11:titration:8}
لن تتفاعل كل قطرة فورًا: فيبقى لون \omterm{def:g11:titration:titration}{المحلول المعايِر}
قبل \omterm{def:g11:titration:equivalence}{التكافؤ}، ويُرى انتهاء \omterm{def:g11:titration:titration}{المعايرة} مبكرًا جدًا (أو يلزم
الانتظار دقائق بعد كل قطرة).
\end{solution}

\begin{solution}{exo:g11:titration:9}
أكبر مما ينبغي: لقد تجاوزت \omterm{def:g11:titration:equivalence}{التكافؤ}، فأضافت البرمنغنات بفائض.
وكان عليها قرب \omterm{def:g11:titration:equivalence}{التكافؤ} أن تضيف \omterm{def:g11:titration:titration}{المحلول المعايِر} قطرةً قطرةً،
مع الرجّ، وأن تتوقف عند أول لون وردي خفيف ثابت.
\end{solution}

\begin{solution}{exo:g11:titration:10}
$n(\ce{Fe^{2+}}) = 5 \times 0.0200 \times 0.0130 = \qty{1.30e-3}{mol}$ في
\qty{10.0}{mL}: $c = \qty{0.130}{mol/L}$ \omterm{def:g2:mixing-and-dissolving:dissolve}{للمحلول} المخفف،
و$10 \times 0.130 = \qty{1.30}{mol/L}$ \omterm{def:g2:mixing-and-dissolving:dissolve}{للمحلول} المركّز.
\end{solution}

\begin{solution}{exo:g11:titration:11}
$0.0630 \times 55.8 = \qty{3.52}{g}$.
\end{solution}

\begin{solution}{exo:g11:titration:12}
$0.1 / 7.2 \approx \qty{1.4}{\percent}$، مقابل
\qty{0.7}{\percent} لحجم \qty{14.3}{mL}. عايِر حجمًا أكبر (أو عيّنة
أشد تركيزًا) من \omterm{def:g2:mixing-and-dissolving:dissolve}{المحلول}، أو استعمل \omterm{def:g2:mixing-and-dissolving:dissolve}{محلولًا} معايِرًا أكثر تخفيفًا.
\end{solution}

\begin{solution}{exo:g11:titration:13}
$n(\ce{MnO4^-}) = 0.0200 \times 0.0176 = \qty{3.52e-4}{mol}$؛
$n(\ce{H2O2}) = \frac{5}{2} \times \num{3.52e-4} = \qty{8.80e-4}{mol}$ في
\qty{10.0}{mL}: \qty{0.0880}{mol/L} في \omterm{def:g2:mixing-and-dissolving:dissolve}{المحلول} المخفف،
و\qty{0.880}{mol/L} في المطهِّر، أي
$0.880 \times 34.0 = \qty{29.9}{g/L}$ (\omterm{def:g2:mixing-and-dissolving:dissolve}{محلول} بنحو 3 \%).
\end{solution}

\begin{solution}{exo:g11:titration:14}
$n(\ce{MnO4^-}) = \num{1.4e-3} / 5 = \qty{2.8e-4}{mol}$، تُصبّ
في نحو \qty{0.015}{L}: $c \approx \qty{0.019}{mol/L}$، إذن
\omterm{def:g2:mixing-and-dissolving:dissolve}{محلول} تركيزه \qty{0.0200}{mol/L}. ولو كان أشد تركيزًا بعشر مرات لكان $V_E$
نحو \qty{1.4}{mL}: ولبلغ خطأ السحاحة قرابة
\qty{7}{\percent}.
\end{solution}

\begin{solution}{exo:g11:titration:15}
$8 \times 0.0200 \times 0.0143 = \qty{2.29e-3}{mol}$ من \omterm{def:g9:ions:ion}{أيونات} الهيدروجين.
و\qty{10}{mL} بتركيز \qty{1.0}{mol/L} تحوي \qty{1.0e-2}{mol}: وهذا يكفي، بل
يزيد على أربعة أضعاف الحاجة. ودون الحمض لا يمكن لتفاعل الفصل أن
يجري كما هو مكتوب (\omterm{def:g9:ions:ion}{فأيونات} الهيدروجين \omterm{def:g7:chemical-reactions:reactant}{متفاعل})، ولتفاعلت البرمنغنات
على نحو آخر.
\end{solution}

\begin{solution}{pb:g11:titration:1}
\textbf{1.} $\ce{MnO4^-}/\ce{Mn^{2+}}$ و$\ce{Fe^{3+}}/\ce{Fe^{2+}}$. \omterm{def:g11:redox:oxidant}{والمؤكسد}
هو \omterm{def:g9:ions:ion}{أيون} البرمنغنات، \omterm{def:g11:redox:oxidant}{والمختزل} \omterm{def:g9:ions:ion}{أيون} الحديد (II).

\textbf{2.} \ce{MnO4^- + 8H+ + 5e- -> Mn^{2+} + 4H2O}؛
\ce{Fe^{2+} -> Fe^{3+} + e-}.

\textbf{3.} تُضرب الثانية في 5:
\ce{MnO4^- + 8H+ + 5Fe^{2+} -> Mn^{2+} + 5Fe^{3+} + 4H2O}. الشحنة على
اليسار $-1 + 8 + 10 = +17$؛ وعلى اليمين $+2 + 15 = +17$.

\textbf{4.} \omterm{def:g9:ions:ion}{أيونات} الهيدروجين \omterm{def:g7:chemical-reactions:reactant}{متفاعل}: فالتفاعل يحتاج إلى \omterm{def:g9:acids-bases-ph:acidic}{محلول
حمضي}.

\textbf{5.} هو تام، وسريع، وهو التفاعل الوحيد للبرمنغنات
في الدورق؛ والبرمنغنات البنفسجية تُستهلك قبل \omterm{def:g11:titration:equivalence}{التكافؤ}
وتلوّن الدورق بالوردي بعده.

\textbf{6.} في السحاحة؛ ويُقرأ عند أسفل الهلالة، والعين على مستواها،
قبل \omterm{def:g11:titration:titration}{المعايرة} وبعدها، بدقة \qty{0.05}{mL}.

\textbf{7.} أصفر باهت (\omterm{def:g9:ions:ion}{أيونات} الحديد)؛ وكل قطرة بنفسجية تستهلكها فورًا
\omterm{def:g9:ions:ion}{أيونات} الحديد (II) التي ما زالت موجودة.

\textbf{8.} $0.0200 \times 0.0143 = \qty{2.86e-4}{mol}$.

\textbf{9.} $n(\ce{Fe^{2+}}) = 5 \times \num{2.86e-4} =
\qty{1.43e-3}{mol}$.

\textbf{10.} $\num{1.43e-3} \times 55.8 = \qty{0.0798}{g} =
\qty{79.8}{mg}$.

\textbf{11.} $(80 - 79.8)/80 \approx \qty{0.3}{\percent}$.

\textbf{12.} $\qty{0.1}{mL}$ من \qty{14.3}{mL} تساوي \qty{0.7}{\percent}،
أي نحو \qty{0.6}{mg}: فالنتيجة، $79.8 \pm 0.6$ mg، تتفق مع
اللصاقة.

\textbf{13.} $14.1$: $5 \times 0.0200 \times 0.0141 \times 55.8 =
\qty{78.7}{mg}$؛ $14.4$: \qty{80.4}{mg}.

\textbf{14.} $(79.8 + 78.7 + 80.4)/3 = \qty{79.6}{mg}$.

\textbf{15.} الأقراص نفسها تختلف قليلًا في محتواها من
الحديد، بنحو \qty{1}{\percent}.

\textbf{16.} لا: سيكون $V_E$ نحو \qty{2.9}{mL}، وخطأ السحاحة
\qty{3.5}{\percent}، أي أسوأ بخمس مرات.

\textbf{17.} سيتأكسد فيتامين C أيضًا بالبرمنغنات: فلا يعود
التفاعل وحيدًا، ويصير $V_E$ أكبر مما ينبغي، وتُقدَّر كتلة
الحديد بأكثر من حقيقتها.

\textbf{18.} $\num{1.43e-3} \times 151.9 = \qty{0.217}{g}$، أي نحو
\qty{217}{mg}.

\textbf{19.} $\num{1.43e-3} \times \num{6.02e23} = \num{8.6e20}$ \omterm{def:g9:ions:ion}{أيون}.

\textbf{20.} نحو \qty{80}{mg} من الحديد في القرص: فاللصاقة
صادقة.
\end{solution}
